% Project Phase 1 report -- IEEE two-column, Times New Roman, 10 pt.

\documentclass[10pt,conference,a4paper]{IEEEtran}
\usepackage[T1]{fontenc}
\usepackage{newtxtext,newtxmath}   % Times New Roman text and math
\usepackage{amsmath}
\usepackage{booktabs}
\usepackage{array}
\usepackage{tabularx}
\usepackage{url}
\usepackage[hidelinks]{hyperref}
\usepackage{caption}

\newcommand{\HB}{\mathrm{HighBits}}
\newcommand{\LB}{\mathrm{LowBits}}
\newcommand{\norm}[1]{\lVert #1 \rVert_\infty}

\begin{document}

\title{Breaking the Threshold: A Systematic Cryptanalysis of Threshold ML-DSA
Submissions to the NIST Multi-Party Threshold Call}

\author{\IEEEauthorblockN{Ashmit Dhown \quad 2023A7PS0205U \qquad Srivathsa H Honyal \quad 2023A7PS0198U}
\IEEEauthorblockA{Project Phase 1 Report -- Cryptography\\
BITS Pilani, Dubai Campus}}

\maketitle

% ---- Contribution statement (required on first page) ----
\begin{table}[h]
\centering
\renewcommand{\arraystretch}{1.3}
\begin{tabular}{|p{3.2cm}|p{4.6cm}|}
\hline
\textbf{Member} & \textbf{Contribution} \\
\hline
Ashmit Dhown\newline 2023A7PS0205U &
  Literature search and reading (Niot, TALUS, Mithril writeups);
  attack surface analysis (Sections~\ref{sec:attacks}, \ref{sec:surface});
  problem formulations P1 and P4;
  report writing and editing. \\
\hline
Srivathsa H Honyal\newline 2023A7PS0198U &
  Literature search and reading (Quorus, Trilithium, Kao);
  attack toolbox survey (Sections~\ref{sec:toolbox}, \ref{sec:compare});
  problem formulations P2 and P3;
  harness design and experimental methodology (Section~\ref{sec:solution}). \\
\hline
\end{tabular}
\caption*{\footnotesize Both members jointly developed the attack checklist (Table~\ref{tab:checklist}),
the literature matrix (Table~\ref{tab:matrix}), and reviewed all drafts.}
\end{table}

% =====================================================================
\begin{abstract}
ML-DSA (FIPS~204) is the post-quantum digital signature standardised by NIST. It was built for one signer holding one
key. Many real systems would rather split the key between several devices, so that stealing one device is not enough to
sign. Doing this for ML-DSA is hard, because its signing process uses rejection sampling and rounding, and these steps do
not work well with secret sharing. In 2026 NIST opened a First Call for multi-party threshold schemes (IR~8214C). Four
proposals in that call produce signatures that a normal ML-DSA verifier accepts: Mithril, Quorus, SplitForge/Trilithium
and TALUS. In this report we survey twenty-four papers and specifications. They cover these four schemes, the work they
build on, and the attack tools (integer learning with errors, integer linear programming and robust regression) that
have been used to recover keys from lattice signatures. We find that early versions of TALUS were broken by two practical
key-recovery attacks. No public attack exists on the other three, but mainly because nobody has tested them
systematically. We state four open problems and propose a reusable testing harness. It simulates what each scheme
publishes, including failed signing attempts, and runs a set of statistical key-recovery tests on that data.
\end{abstract}

\begin{IEEEkeywords}
ML-DSA, FIPS 204, threshold signatures, lattice cryptography, cryptanalysis, integer LWE, rejection sampling.
\end{IEEEkeywords}

% =====================================================================
\section{Introduction}

\textbf{Digital signatures and quantum computers.} A digital signature proves that a message came from a particular
signer and was not changed on the way. The signatures used today, RSA and ECDSA, would be broken by a large quantum
computer. For this reason NIST has standardised new post-quantum signatures. The main general-purpose one is ML-DSA,
defined in FIPS~204~\cite{fips204} and based on CRYSTALS-Dilithium~\cite{dilithium}. Its security comes from two lattice
problems called Module-LWE (learning with errors) and Module-SIS.

\textbf{How ML-DSA signs.} The public key contains a matrix $A$ of size $k\times\ell$. Its entries are polynomials in the
ring $R_q=\mathbb{Z}_q[X]/(X^{256}+1)$ with $q=8380417$. The secret key is a pair of short vectors $(s_1,s_2)$, and the
public key is $t=As_1+s_2$ with its low bits removed. To sign a message, the signer picks a random short vector $y$ (the
nonce) and computes $w=Ay$. It keeps the top part $w_1=\HB(w)$ and hashes it with the message to get a challenge $c$, which
is a sparse polynomial. The response is $z=y+cs_1$, sent together with a small hint $h$.

Before sending, the signer runs two checks. The first, $\norm{z}<\gamma_1-\beta$, makes sure $z$ does not reveal $s_1$. The
second, $\norm{\LB(w-cs_2)}<\gamma_2-\beta$, protects $s_2$ and keeps the hint small. If either check fails, the signer
throws the attempt away and starts again with a new nonce. This is called ``Fiat--Shamir with aborts''. Its purpose is
simple. Every $z$ that passes the check looks the same no matter what the secret key is, so signatures do not leak the key.

\textbf{Parameters.} The three ML-DSA versions all use $q=8380417$, $n=256$ and drop $d=13$ bits of $t$. The other values are
listed in Table~\ref{tab:params}. Here $\tau$ is the number of nonzero ($\pm1$) entries in $c$, and $\beta=\tau\eta$ is the
largest size that $cs$ can have.

\begin{table}[t]
\caption{ML-DSA parameter sets~\cite{fips204}}\label{tab:params}
\centering\footnotesize
\begin{tabular}{@{}lcccccc@{}}
\toprule
Set & Cat. & $(k,\ell)$ & $\eta$ & $\tau$ & $\gamma_1$ & $\gamma_2$\\
\midrule
ML-DSA-44 & 2 & (4,4) & 2 & 39 & $2^{17}$ & $(q-1)/88$\\
ML-DSA-65 & 3 & (6,5) & 4 & 49 & $2^{19}$ & $(q-1)/32$\\
ML-DSA-87 & 5 & (8,7) & 2 & 60 & $2^{19}$ & $(q-1)/32$\\
\bottomrule
\end{tabular}
\end{table}

\textbf{Threshold signatures.} In a $(T,N)$ threshold scheme the key is split among $N$ parties. Any $T$ of them can sign
together, but a group smaller than $T$ learns nothing about the key. The classic tool is Shamir secret
sharing~\cite{shamir}. The secret is hidden as the constant term of a random polynomial of degree $T-1$, and each party
gets one point on it. Feldman~\cite{feldman} added public commitments $g^{a_i}$ so that each party can check its share.
For Schnorr signatures, schemes such as FROST~\cite{frost} work well because signing is linear, so partial signatures
just add up.

ML-DSA does not have these helpful properties. A share of a short secret is usually not short, and short values are
exactly what rejection sampling needs. $\HB$ is not linear, so each party cannot round its own piece of $w$. The second
check needs the whole secret $s_2$. Finally, as Section~\ref{sec:attacks} shows, the lattice version of a Feldman
commitment, $A\cdot x$, does not hide $x$ at all.

\textbf{Why look for attacks now.} NIST's First Call (IR~8214C)~\cite{ir8214c} invites threshold schemes. Preliminary
packages are due from late 2026, and full packages not before March 2027. Finding problems is most useful before the
designs are fixed. Two simple facts make ML-DSA threshold designs fragile. First, $A$ is ``tall'' ($k\ge\ell$ in every
version). So if anyone publishes $Ax$ for a secret $x$, then $x$ can be found by ordinary Gaussian elimination, with no
hard lattice problem involved. Second, the equation $z=y+cs_1$ holds over the integers. So if the nonce $y$ ever leaks,
the key follows directly as $s_1=c^{-1}(z-y)$.

\textbf{Contributions.} This report gives (i) an organised survey of threshold ML-DSA designs and of the tools used to
attack them, (ii) a list of gaps in current knowledge, (iii) four clearly stated open problems, and (iv) a proposed
solution, a leakage-testing harness, together with a first small experiment.

\textbf{How we chose the papers.} We searched in three rounds. First we took the four proposals and their preview
writeups from the NIST First Call page. Next we followed each proposal's ePrint or arXiv paper and the earlier work it
cites. Finally we searched ePrint, TCHES, EUROCRYPT and ASIACRYPT for terms such as ``Dilithium rejected signatures
ILP'', ``ILWE'' and ``threshold lattice signatures''. We kept a source if it builds a threshold ML-DSA scheme, attacks one,
or gives a technique that such an attack uses. We checked the details of 22 of the 24 references online. The two classic
papers (Shamir and Feldman) are cited from memory. Most sources are from 2024 to 2026, because the topic hardly existed
before that.

\textbf{Scope.} This is a study of papers and specifications. We did not run any scheme's own code. Our only experiment is
a small simulation of the least-squares attack on made-up TALUS-style data (Section~\ref{sec:prelim}). Where we describe
possible weak points, we mark them as hypotheses.

% =====================================================================
\section{Literature Survey}

\subsection{Background: lattice threshold signatures}
Dilithium~\cite{dilithium} introduced the signing method described above, and FIPS~204~\cite{fips204} turned it into
ML-DSA. The first lattice threshold signatures used heavy machinery. G\"ur, Katz and Silde~\cite{gks}, for example, build a
two-round scheme from threshold homomorphic encryption. Threshold Raccoon~\cite{raccoon} (EUROCRYPT 2024) is much
lighter. It uses only symmetric primitives and simple lattice operations, and adds one-time masks so that partial
signatures do not leak partial keys. Its signatures, however, are not ML-DSA signatures. Del Pino and
Niot~\cite{finally} give a compact Raccoon-based scheme, limited to $2^{64}$ signatures and about 8 parties. Del Pino et
al.~\cite{dpia} add identifiable aborts, but that paper has since been withdrawn in favour of a later one. The important
recent change is that new schemes aim to produce signatures that the \emph{normal} FIPS~204 verifier accepts. This
matters because certificates, hardware and pinned keys that are already deployed cannot easily be changed.

\textbf{Rounding steps that matter for attacks.} ML-DSA shortens the public key with Power2Round: $t=t_1 2^d+t_0$, and only
$t_1$ is published. While signing, Decompose splits $w$ into $w_1$ and a low part that is at most $\gamma_2$ in size. The
verifier computes $v=Az-ct_1 2^d$, which equals $w-c(s_2-t_0)$. So $v$ is public and contains the term $c(s_2-t_0)$. The
second rejection check keeps $\LB(w)$ at least $\beta$ away from a rounding boundary, and this is what stops $v$ from
giving away information. If the check is removed or changed, that protection is gone.

\subsection{The four First-Call ML-DSA proposals}

\textbf{Mithril}~\cite{mithril,reloaded,mithrilposter,mithrilpw} (PQShield, Brave, Bristol) was the first practical
threshold scheme compatible with ML-DSA. It uses \emph{replicated} secret sharing, where shares stay short, so they fit
ML-DSA's size limits. Each party does its own rejection check using samples from a hyperball, so no expensive joint
computation is needed for aborts. It needs about 1\,MB of communication per party for up to six parties and comes with a
Go implementation. Its security proof reduces to MLWE and to the unforgeability of ML-DSA, for an attacker who corrupts
up to $T-1$ parties chosen in advance. There are no identifiable aborts. Each attempt succeeds with probability $1/2$
rather than about $1/4$ in plain ML-DSA, so the distribution of accepted responses is not the same as in FIPS~204. Each
party publishes a full MLWE sample as its commitment, so the simple $Ax$ attack does not apply. The proof uses R\'enyi
divergence and only covers $Q_s=2^{50}$ signatures. The paper appears at USENIX Security 2026.

\textbf{Quorus}~\cite{quorus} (J.P.\ Morgan) changes the signing process so that it is easier to run inside multi-party
computation (MPC), while keeping normal verification and signature size. It assumes an honest majority (fewer than $n/2$
corrupt parties) and has been tested with up to 63 parties. Online cost is about 150\,KB per party per attempt, or 0.31 to
0.59\,MB per final signature, over 16 or 29 rounds. Security is proven in the universal composability (UC) framework. In
its variant the nonce is drawn uniformly, a small noise term $e_w$ is added to $w$, and $w_1$ is always output, even when
$(z,h)$ is rejected. The proof argues that these rejected outputs can be simulated and so reveal nothing.

\textbf{SplitForge/Trilithium}~\cite{trilithium} (Cybernetica) is for exactly two parties, a server and a phone, plus a
helper that supplies correlated randomness (the CRP). It is UC-secure if at most one of the three is malicious. It uses a
secure comparison protocol and a new rounding protocol, and has a Rust implementation. Key generation takes 3 rounds and
each signing attempt 14. It publishes the high bits $w_H$ even when an attempt is rejected. To argue this is safe it relies
on an MLWR-type assumption, which the paper itself calls non-standard for ML-DSA parameters.

\textbf{TALUS}~\cite{talus,taluspw} (Codebat) introduces the \emph{Boundary Clearance Condition} (BCC). For a fixed share of
nonces (31.7\% at ML-DSA-65, 43.2\% at ML-DSA-44 and 39.1\% at ML-DSA-87) it can be proven that $s_2$ cannot push $w$
across a rounding boundary. So the rejection checks can be done ahead of time on prepared nonces. This gives a
distributed MPC version with only two online rounds. A one-round version using trusted hardware is described in v0.22
only as a comparison. The paper says its protection of $s_2$, a cap on the number of signatures, is supported by a model
rather than proven. It also proves a useful lower bound: any FIPS-exact scheme that reveals a nonce made by adding up
shares (an Irwin--Hall nonce) can have its key recovered after about $2^{30}$ signatures. TALUS builds on Kao's Shamir
nonce DKG~\cite{kao}. There the nonce itself is Shamir-shared and hidden by masks that cancel in pairs, and the paper
claims the nonce shares have more than five times the key's entropy for signing groups of up to 17.

\textbf{Why local rejection is tricky in Mithril.} In normal ML-DSA a nonce is accepted only if both checks pass, which
happens about $1/4$ of the time. With $N$ parties, each holding a short share, the \emph{sum} of their responses has to land
in the same box, but no single party knows that sum. Mithril's answer is to let each party draw its nonce share from a
hyperball (a round ball instead of a box). The size of the total is then predictable, and each party can decide locally.
The catch is that a sum of ball-shaped pieces is not uniform on a box. The accepted responses are therefore close to the
FIPS~204 distribution, but not equal to it. How big this difference is, and whether an attacker could measure it with a
realistic number of signatures, is our problem P2. The difference is part of the design, not a bug, so the question is
``how much'', not ``yes or no''.

Table~\ref{tab:schemes} compares the four proposals side by side.

\begin{table*}[t]
\caption{The four ML-DSA-compatible First-Call proposals}\label{tab:schemes}
\centering\footnotesize
\begin{tabularx}{\textwidth}{@{}lXXlX@{}}
\toprule
Scheme & Design & Trust / corruption model & Rounds & Public attacks\\
\midrule
Mithril~\cite{mithril} & Replicated short shares, local rejection & Dishonest majority ($T-1$), static; $N\le6$ & $\sim$3 per attempt & None found\\
Quorus~\cite{quorus} & MPC-friendly ML-DSA variant & Honest majority, UC; up to 63 parties & 16 or 29 & None found\\
Trilithium~\cite{trilithium} & 2-party MPC + correlated randomness & CRP trusted; one party malicious & 14 per attempt & None; relies on heuristic assumption\\
TALUS~\cite{talus,taluspw} & BCC nonce filtering, Shamir nonce DKG & Honest majority, $N\ge 2T-1$ & 2 (v0.22) & Earlier versions broken~\cite{niot}\\
\bottomrule
\end{tabularx}
\end{table*}

\subsection{Published attacks on TALUS}\label{sec:attacks}
Niot~\cite{niot} (ePrint 2026/1386) describes two separate, practical attacks on early versions of TALUS.

\textbf{Attack 1: commitments that do not hide.} TALUS-MPC used Feldman-style commitments of the form $A\cdot x$, an idea
copied from discrete-log protocols. But $A$ can be inverted from the left, except with probability at most $2^{-15}$
(ML-DSA-44) or $2^{-38}$ (ML-DSA-65 and 87). So $x=A^{+}(Ax)$. The commitment $As_{1,i}$ sent during key generation
therefore gives away every key share. In the same way, $A\hat y$ sent in the blame phase gives away the nonce, and then
$s_1=c^{-1}(z-y)$ follows from just \emph{one} signature.

\textbf{Attack 2: the missing $s_2$ check.} Without the second check, each signature releases $v=w-c(s_2-t_0)$. The hint
reveals $\HB(w)$, so what the attacker really sees is $b=-cs'+e$, where $s'=s_2-t_0$ and $e=\LB(w)$ is close to uniform
noise. This is an ILWE problem (LWE without the ``mod $q$''), and Bootle et al.~\cite{ilwe} showed it can be solved by least
squares. The number of signatures needed is about
\begin{equation}
N \gtrsim \frac{4\gamma_2^2}{3\tau}\label{eq:ilwe}
\end{equation}
which is roughly $3.1\times10^8$ for ML-DSA-44, $1.9\times10^9$ for ML-DSA-65 and $1.5\times10^9$ for ML-DSA-87. Niot points
out that these numbers are not optimised. He also notes that simply putting the check back does not fix the problem, since
a corrupt party can choose a biased nonce share (for example $y_i=0$) so that the check always passes.

The TALUS authors responded in preview writeup v0.22~\cite{taluspw}. They replaced the commitments with zero-knowledge
proofs that shares are well formed. For the $s_2$ problem they added a hard limit on signatures per key, about
$2^{13}$, $2^{14}$ and $2^{14}$ for the three versions, enforced by changing keys. They compare this limit with a point they
call the ``integer-uniqueness wall'', near $2^{15.2}$, $2^{16.4}$ and $2^{16.2}$. Each prepared nonce is tied to one group of
signers and deleted after use. Whether the safety margin is real is exactly the kind of question this project asks. The
limit also makes the scheme harder to deploy, because certificates and pinned keys cannot be replaced every few thousand
signatures.

\textbf{Attack 1 in more detail.} Suppose party $i$ sends $C_i=As_{1,i}$. Using the NTT, $A$ breaks into 256 small
$k\times\ell$ matrices over $\mathbb{Z}_q$, one for each slot. Each one is invertible unless something with probability
about $1/q$ happens, which is where the $2^{-15}$ and $2^{-38}$ numbers come from. Anyone watching can solve each small
system and get $s_{1,i}$ exactly. In the simplest case, $k=\ell=1$, the commitment is just $a\cdot s \bmod q$, and dividing
by $a$ gives back $s$. This shows why the idea ``discrete-log commitments hide, so lattice ones should too'' is wrong.
$g^x$ hides $x$ because discrete logarithms are hard. $Ax$ is just a linear map, and linear maps are easy to undo.

\textbf{Attack 2 in more detail.} Look at one coefficient of $s'$. Each signature gives $b=-cs'+e$, where $c$ has a few
$\pm1$ entries and $e$ is bounded noise. One sample only says that a sum of $\tau$ secret coefficients lies in a window about
$2\gamma_2$ wide. That window is much wider than the range of the secret, so a single sample tells us almost nothing.
But when we average many samples with different challenges, the uncertainty shrinks like $1/\sqrt{N\tau}$. Once it is
less than half the gap between whole numbers, each coefficient can be rounded to the right value. The ratio
$\gamma_2^2/\tau$ in (\ref{eq:ilwe}) also explains why ML-DSA-44, whose $\gamma_2=(q-1)/88$ is small, is more exposed per
signature than the other two versions.

\textbf{Why adding up nonces leaks.} In many designs each of the $T$ parties adds a uniform random share to the nonce. The
total is not uniform. It is bell-shaped (an Irwin--Hall distribution). Then $z=y+cs_1$ is that bell shape moved by $cs_1$,
so each signature gives away a tiny amount of information about $s_1$. In~\cite{kao} this is less than 0.013 bits per
signature, but it adds up, and~\cite{talus} proves that an efficient attack works after about $2^{30}$ signatures at
ML-DSA-65. A design can avoid this in only two ways. It can make the revealed nonce exactly uniform, which is what Quorus
and Trilithium claim to do. Or it can limit the number of signatures, which is what the TALUS cap does.

\textbf{Who the attacker is.} Proofs make different assumptions about the attacker. With \emph{static} corruption the
attacker picks which parties to control before the protocol starts. With \emph{adaptive} corruption it can pick during
the run. Dishonest-majority security (Mithril) holds even if up to $T-1$ signers are corrupt. Honest-majority security
(Quorus, TALUS-MPC) needs $N\ge2T-1$ and promises nothing beyond that. Identifiable aborts let honest parties name a
cheater. Without them, as in Mithril, a corrupt party can see the honest parties' messages, abort, and try again, so the
data from aborted attempts is also in the attacker's hands. Each of these choices changes what an attack has to look at,
so our harness takes the attacker model as an input.

\textbf{Checking the numbers.} Putting the parameters into (\ref{eq:ilwe}) gives the same figures as~\cite{niot} (see
Table~\ref{tab:ilwe}). That is about $2^{28}$ to $2^{31}$ signatures, or $2^{12}$ to $2^{15}$ times the TALUS cap. Two points
are worth stressing. The estimate treats the noise as if it were Gaussian, so it ignores the useful information in the
noise's sharp edges. It also finds $s'$ one coefficient at a time, while lattice reduction can use the fact that all
coefficients are small at once. Whether a better attack can close this 12 to 15 bit gap is our problem P1.

\subsection{Attacks on lattice signatures through leakage}\label{sec:toolbox}
The tools behind these attacks come from earlier work on ILWE and on leaks from rejection sampling. Bootle et
al.~\cite{ilwe} first framed linear leakage in the BLISS signature as an ILWE problem. Zhou et al.~\cite{rejected} (TCHES
2025) showed that the challenges of \emph{rejected} Dilithium signatures leak. Each rejection gives upper and lower bounds
on $cs$. They turn these bounds into an integer linear program (ILP) and recover full keys in seconds to minutes, which
they confirmed on an ARM Cortex-M4 chip. Damm et al.~\cite{concealed} (ASIACRYPT 2025) studied ``Concealed ILWE'', where
only some of the samples carry information. They showed that ordinary least squares fails there, and used Huber and
Cauchy robust regression to break a masked Dilithium in under two minutes. Yates et al.~\cite{yates} tried least-squares
attacks on ILWE built from real signatures and found that, at the parameters they studied, the schemes held up. Taken
together, these papers show that any threshold design that releases values related to $cs_1$ or $cs_2$ should be tested
with stronger tools than plain least squares.

\textbf{Defences.} Masking and shuffling protect against side channels, but~\cite{concealed} shows that masking alone is
not enough if some samples still leak. Against leaks from rejections, the defence is to hide which check failed and to
release nothing computed from a rejected response. This is why Quorus, which releases $(w_1,c,\bot)$, needs a careful
argument that this is safe. Each defence has a cost: masking costs time, a signature cap costs key changes, and hiding
rejected attempts costs extra communication.

\subsection{Comparing the proposals}\label{sec:compare}
Each proposal makes a different trade-off between how many corrupt parties it tolerates, how many parties it supports,
how much communication it needs and how strong its evidence of security is. Only Mithril tolerates a dishonest majority,
but it supports at most six parties. Quorus scales to many parties and has UC proofs, but needs an honest majority and
many rounds. Trilithium covers the phone-and-server case but needs an extra trusted helper. TALUS has the fewest rounds,
but its security is the least settled, since the $s_2$ problem is handled by a usage limit rather than by the design
itself. The quality of the evidence also differs. Mithril and Quorus have security reductions. Trilithium has a UC proof
that rests on a heuristic. TALUS has an information-theoretic argument with a numerical margin. Independent testing
therefore helps most where a proof relies on a heuristic or a margin (Trilithium, TALUS), and where the parameters change
the accepted distribution (Mithril).

\subsection{Related and older work}
Mithril started as a poster~\cite{mithrilposter} and an early ePrint~\cite{reloaded}. That ePrint covered both ML-DSA
and an improved Raccoon with identifiable aborts (up to 64 parties), and was later split into two papers, one of which
is~\cite{mithril}. The preview writeup~\cite{mithrilpw} adds that up to 8 parties is still practical, that no
synchronised broadcast channel is needed, and that the Go code uses floating-point numbers for hyperball sampling.

Kao's earlier paper~\cite{kao} is not a Call submission, but it shows the same patterns. Its fully distributed profile P2
broadcasts masked commitments. The paper's own remarks show that if $\lambda_h A y_h$ is ever exposed, it gives away $y_h$
and then the key. The masks only hide these values when $|S\setminus C|\ge2$, that is, when at least two signers are
honest. The paper admits that when $T=N$ and $N-1$ parties are corrupt, the masks do not hide anything in P2. Yet its
Table~1 still claims dishonest-majority unforgeability with UC proofs against static attackers. Its Irwin--Hall analysis
limits the loss to 0.013 bits per signature, and says the bound is meaningful for fewer than 16{,}000 signatures. Later,
\cite{talus} proved a passive attack after about $2^{30}$ signatures for any scheme that reveals such a summed nonce.

\subsection{Possible weak points of each scheme}\label{sec:surface}
This section lists where we plan to look first. These are hypotheses based on the specifications, not results.

\textbf{Mithril.} (i) The accepted responses: since each attempt succeeds with probability $1/2$, the information leaked
must be measured beyond the proven $2^{50}$ signatures, up to the usual $2^{64}$. (ii) Key sharing after the fact gives
the attacker a noisy hint about the ML-DSA secret. The authors put this at a 7 to 12 bit loss. This should be checked
for the exact shares that $T-1$ corrupt parties hold. (iii) Adaptive corruption is only argued informally, with a loss of
at most 5 bits for $N\le6$. This is a gap in the proof, not an attack. (iv) Without identifiable aborts, corrupt parties
can abort after seeing the honest parties' messages. (v) The floating-point hyperball sampler could be slightly biased or
leak through timing. A fixed-point C version is planned.

\textbf{Quorus.} (i) Rejected attempts release $(w_1,c,\bot)$. The paper proves this is safe, and we can test it in
practice with regression and ILP against $cs_1$ and $cs_2$. (ii) Zhou et al.~\cite{rejected} show that if an attacker
learns \emph{which} check failed, the key can be recovered, so implementations must reveal only $\bot$. (iii) The MPC
building blocks (the joint rejection check, batched OR and offline preparation) are where coding mistakes would hide.
(iv) The edge case $N=2T-1$, where an honest majority just barely holds, deserves a check.

\textbf{SplitForge/Trilithium.} The paper publishes $w_H$ even for rejected attempts. It openly says this needs an
MLWR-type assumption, and notes that other work either adds noise (like Quorus), adds a new assumption, or leaves the
question open. We think this is the claim most worth attacking. Since $y$ is new in every attempt, each $w_H$ is an
LWR-style sample under a fresh secret. Any build-up of information about the key must therefore come from the
rejection events. The rejection result is revealed to both parties, so the information it carries about $cs_2$ should be
bounded per attempt. The CRP is a trusted third party. According to the writeup it can carry out selective-disclosure
attacks, which the authors consider harmless.

\textbf{TALUS v0.22.} The signature cap is the whole argument for protecting $s_2$. Under BCC the noise $e=\LB(w)$ lies in
$(-\gamma_2+\beta,\gamma_2-\beta)$. Noise with sharp edges like this gives away more per sample than its variance
suggests. Other places to look are the zero-knowledge proofs, the blame procedure, the binding of nonces to one signing
group (a nonce-reuse issue was reported on an earlier version), and corrupt parties biasing their nonces.

\subsection{Literature matrix}
Table~\ref{tab:matrix} rates every source by relevance (H central, M supporting, L background) and marks which themes it
covers: TS threshold signatures, MD ML-DSA compatible, AT attack or leakage, SS secret sharing or MPC, RS rejection
sampling. Sixteen of the 24 sources are from 2025 or 2026, and 22 are from 2018 or later.

\begin{table}[t]
\caption{Literature matrix}\label{tab:matrix}
\centering\scriptsize
\begin{tabular}{@{}llcccccc@{}}
\toprule
Ref & Year & Rel. & TS & MD & AT & SS & RS\\
\midrule
\cite{fips204} FIPS 204 & 2024 & H &  & x &  & x & \\
\cite{dilithium} Dilithium & 2018 & M &  & x &  & x & \\
\cite{ir8214c} IR 8214C & 2026 & H & x & x &  &  & \\
\cite{mithril} Mithril & 2026 & H & x & x &  & x & x\\
\cite{reloaded} TS Reloaded & 2025 & M & x & x &  &  & x\\
\cite{mithrilposter} Mithril poster & 2025 & L & x & x &  &  & \\
\cite{mithrilpw} Mithril PW & 2026 & H & x & x &  & x & x\\
\cite{quorus} Quorus & 2025 & H & x & x &  & x & x\\
\cite{trilithium} Trilithium & 2025 & H & x & x &  & x & x\\
\cite{talus} TALUS & 2026 & H & x & x & x & x & x\\
\cite{taluspw} TALUS PW v0.22 & 2026 & H & x & x & x & x & x\\
\cite{kao} Shamir nonce DKG & 2026 & H & x & x & x & x & \\
\cite{niot} Niot attack & 2026 & H & x & x & x &  & x\\
\cite{raccoon} Thr.\ Raccoon & 2024 & M & x &  &  & x & \\
\cite{dpia} Lattice TS+IA & 2025 & L & x &  &  & x & \\
\cite{gks} G\"ur et al. & 2024 & L & x &  &  & x & \\
\cite{ilwe} ILWE/BLISS & 2018 & H &  &  & x &  & \\
\cite{rejected} Rejected sigs & 2025 & H &  & x & x &  & x\\
\cite{concealed} Concealed ILWE & 2025 & H &  & x & x &  & x\\
\cite{yates} ILWE+rej.\ samp. & 2025 & M &  & x & x &  & x\\
\cite{shamir} Shamir & 1979 & L &  &  &  & x & \\
\cite{feldman} Feldman & 1987 & M &  &  &  & x & \\
\cite{frost} FROST & 2020 & L & x &  &  & x & \\
\cite{finally} Finally! & 2025 & M & x &  &  & x & \\
\bottomrule
\end{tabular}
\end{table}

Of the 24 sources, 13 are central, 6 supporting and 5 background. The sources marked AT are the basis for the attack
methods in Section~\ref{sec:solution}.

\textbf{Lessons from the survey.} The failures we found lead to four simple rules. First, never publish $A$ times a secret
or a nonce, however it is masked, unless the mask still hides it in the worst case of corruption allowed. Second, treat
the nonce as being as secret as the key, including during blame and abort steps. Third, keep every rejection check that
the single-signer scheme has, or prove that whatever replaces it leaks no more. Fourth, if the number of signatures per
key is limited, the limit needs a margin that a strong attack, not just least squares, cannot close. None of these rules
is new on its own. But the First-Call submissions show how easy it is to break them when ideas are copied from
discrete-log threshold schemes.

\subsection{Gaps in current knowledge}
\begin{itemize}
\item \textbf{No public testing} of what Mithril, Quorus or Trilithium actually publish. There are only arguments about
the designs.
\item \textbf{TALUS's cap} is about 2 bits below the authors' own wall and 14 to 17 bits below Niot's unoptimised estimate.
A better bounded-noise or ILP attack might close this gap.
\item \textbf{Data from rejected attempts} (Quorus releasing $w_1$ on rejection, Trilithium's MLWR-type assumption) has not
been tested for leaks in any paper we found.
\item \textbf{Mithril's accepted responses} and its key-sharing hint loss (reported as 7 to 12 bits) have not been checked
independently. The proof covers $2^{50}$ signatures, fewer than the usual $2^{64}$.
\item \textbf{Gaps in the attacker model}: proofs that only cover static corruption, the honest-majority edge case
($N=2T-1$), and side conditions such as $|S\setminus C|\ge2$ in Kao's profile P2~\cite{kao}.
\end{itemize}

% =====================================================================
\section{Problem Formulation}

\subsection{Attacker model and notation}
Fix one ML-DSA parameter set $(q,k,\ell,\tau,\gamma_1,\gamma_2,\beta)$. A $(T,N)$ threshold scheme $\Pi$ produces a
\emph{transcript}: everything an attacker can see while a message $m$ is signed. This includes the final signature, all
broadcasts and coordinator messages, the internal state of corrupt parties and, importantly, the data from
\emph{aborted} attempts. We write its distribution as $\mathrm{Tr}_\Pi(\mathit{sk},m)$. The attacker corrupts as many parties
as the scheme allows, chosen in advance. It watches $q_s$ signing sessions on messages of its choice, and wins if it
recovers $s_1$, which is the same as recovering the secret key. We focus on full key recovery, because it is much
stronger evidence of a break than a small statistical bias. Table~\ref{tab:notation} lists the notation.

\begin{table}[t]
\caption{Notation}\label{tab:notation}
\centering\footnotesize
\begin{tabularx}{\columnwidth}{@{}lX@{}}
\toprule
Symbol & Meaning\\
\midrule
$A,t$ & Public $k\times\ell$ matrix over $R_q$; public key $t=As_1+s_2$\\
$s_1,s_2$ & Short secrets; $s'=s_2-t_0$, $t_0$ the dropped low bits of $t$\\
$y,z,c$ & Nonce, response $z=y+cs_1$, sparse challenge ($\tau$ entries $\pm1$)\\
$w,w_1,e$ & $w=Ay$, $w_1=\HB(w)$, rounding remainder $e=\LB(w)$\\
$(T,N)$ & Threshold and parties; $S$ signing set, $C$ corrupt set\\
$q_s$ & Number of signing sessions observed\\
\bottomrule
\end{tabularx}
\end{table}

The schemes differ most in what they publish. Table~\ref{tab:transcript} shows, as we read the sources, what each one
releases apart from the final signature $(c,z,h)$.

\begin{table}[t]
\caption{What each scheme publishes beyond the signature}\label{tab:transcript}
\centering\footnotesize
\begin{tabularx}{\columnwidth}{@{}lXX@{}}
\toprule
Scheme & Published & Aborted attempts\\
\midrule
Mithril & Round messages; per-party commitments to $w$ & Visible to corrupt parties; no identifiable aborts\\
Quorus & MPC messages; honest-majority broadcasts & Releases $(w_1,c,\bot)$\\
Trilithium & Two-party messages; CRP values & Rejection bit revealed; partial responses\\
TALUS v0.22 & ZK well-formedness proofs; blame data & Nonces filtered offline; tied to one group\\
\bottomrule
\end{tabularx}
\end{table}

This table guides the experiments in two ways. The right column tells us whether testing rejected attempts makes sense at
all. TALUS filters nonces offline, so for TALUS the leak to study is in accepted signatures and the cap. The middle column
tells us which checklist items apply. For example, item~1 (published $Ax$) matters for any scheme that broadcasts
commitments.

\subsection{Problem statement}
\emph{For each ML-DSA-compatible threshold scheme in the First Call, and using only its public specification, decide
whether what it publishes (including aborted attempts and what a group of corrupt parties sees) reveals enough about the
secret key to recover it with a realistic number of signatures. If it does not, measure the safety margin in signatures
and in bits of security.} We break this into four open problems.

\textbf{P1 (TALUS signature cap).} Below a point the authors call the integer-uniqueness wall
($q_\mathrm{uniq}\approx2^{15.2}$, $2^{16.4}$, $2^{16.2}$), the noisy values $O=r_0+c(t_0-s_2)$ are not enough on their own to
pin down $s_2$. The cap of $2^{13}$ to $2^{14}$ is about two bits below this point. But the public key already determines
$s_2$, so the real protection is computational. The authors model a meet-in-the-middle attack costing $2^{H/2}$, with
remaining entropy $H\ge502$. The only proven lower bound, however, is about $2^{58}$ at ML-DSA-65 and $2^{49}$ at ML-DSA-87,
which is below 128 bits. For ML-DSA-44 ($k=\ell$) there is no such bound at all. P1 asks for the true cost of recovering
$s'$ from at most $2^{14}$ observations with sharp-edged noise, using bounded-noise methods, ILP and lattice estimators
that take hints into account.

\textbf{P2 (leaks from Mithril's accepted responses).} Compute the Fisher information $I$ about $s_1$ in one accepted
response, given Mithril's hyperball sampling and $1/2$ acceptance rate. From it, find the attack point, about $4/(I\tau)$
signatures, and compare it with the proven $Q_s=2^{50}$ and the usual $2^{64}$. Also recompute how much security is lost
through the key-sharing hint, for the exact shares that $T-1$ corrupt parties hold.

\textbf{P3 (leaks from rejected attempts).} Quorus and Trilithium release $w_H$ and $c$ for rejected attempts. Mithril
releases each party's commitment even when that party's response is rejected. Quorus adds noise $e_w$ so that this can
be proven safe. Trilithium keeps ML-DSA unchanged and relies on an MLWR-type assumption. P3 asks whether the data from
rejected attempts, together with the accept/reject bit, depends on $s_1$ or $s_2$ in any measurable way. It also asks how
many bits about $cs_2$ Trilithium's revealed rejection bit leaks per attempt.

\textbf{P4 (Kao's profile P2).} The paper~\cite{kao} says that in profile P2 the broadcast commitments stay hidden only if
$|S\setminus C|\ge2$. Its own remarks show that an exposed $\lambda_h A y_h$ gives $y_h$, and then
$s_{1,h}=c^{-1}(z_h-y_h)$. If $|S|=T$ and $T-1$ of the signers are corrupt, then $|S\setminus C|=1$. Yet the paper claims
security against a dishonest majority. P4 asks whether P2 is secure only when $|S|\ge T+1$.

\subsection{What we expect and how we will judge success}
For P1 we expect attacks that use hints to do much better than the least-squares estimate. Success means a measured cost
for recovering $s'$ from $2^{14}$ observations, checked against the public key at reduced size. For P2 we expect the attack
point to be above $2^{64}$, so success means a numerical bound with a clear derivation. For P3 we expect Quorus to be safe
as claimed, while Trilithium's assumption may only hold roughly. Success means a statistical test with a stated false
alarm rate and sample budget. For P4 we expect P2 to fail when $|S\setminus C|=1$. Success means either an actual key
recovery on a small example or a clear reason why the masks still work. Writing these down now stops the project from
drifting towards whatever is easiest to measure.

\subsection{Research questions answered by the survey}
(RQ1) Which threshold ML-DSA designs exist, and how do they differ? There are four First-Call proposals. They differ in
attacker model, number of parties and rounds (Table~\ref{tab:schemes}). (RQ2) Which are known to be broken? Early
versions of TALUS, by two separate attacks~\cite{niot}. No public attacks exist on the others. (RQ3) Which techniques can
break such a scheme? Inverting a published $Ax$, recovering a leaked nonce, ILWE by least squares, ILP on rejection
bounds, and robust regression~\cite{niot,ilwe,rejected,concealed}. (RQ4) Where is the evidence weakest? Where proofs use
heuristics (Trilithium), margins (the TALUS cap) or a changed distribution (Mithril). These answers lead directly to P1
to P4.

\textbf{How this differs from earlier attacks.} Earlier attacks on lattice
signatures~\cite{ilwe,rejected,concealed,yates} target one signer whose device leaks through a physical side channel, such
as power use. Here the leak is \emph{logical}. It is part of the protocol messages, so anyone watching, or any corrupt
party, can see it, whatever the hardware. The attacker is a protocol participant, and the noise is known exactly from the
specification instead of being measured. So results apply to every implementation, and schemes that so far have only a
proof get an independent practical check.

% =====================================================================
\section{Possible Solution: A Leakage-Testing Harness}\label{sec:solution}

\subsection{Design}
We propose a reusable piece of software, a ``harness'', with three layers.
(1)~\textbf{ML-DSA building blocks}: parameters, NTT over $R_q$, Power2Round, Decompose, MakeHint/UseHint and the FIPS~204
samplers, checked against the official test vectors.
(2)~\textbf{Scheme simulators}: each one produces the public transcript of a scheme on made-up keys. This covers accepted
and aborted attempts, and what a group of corrupt parties would see.
(3)~\textbf{Attack methods}, one for each item of the checklist in Table~\ref{tab:checklist}: Gaussian elimination for a
published $Ax$, direct recovery from a leaked nonce, least squares (ILWE), bounded-noise and ILP methods, robust
(Huber/Cauchy) regression~\cite{concealed}, Fisher-information calculation, and lattice-estimator accounting for hints.
Each attack method takes a stream of transcript records and returns a guess for the secret with a confidence score. So
adding a new scheme only needs a new simulator. Because the keys are made up, we always know the right answer, and every
claimed recovery is checked by recomputing $t$ from the guessed $s_1$.

\begin{table}[t]
\caption{Attack checklist from the survey}\label{tab:checklist}
\centering\footnotesize
\begin{tabularx}{\columnwidth}{@{}lXX@{}}
\toprule
\# & Question & Typical result\\
\midrule
1 & Is $Ax$ published for a secret or nonce $x$? & $x=A^{+}(Ax)$\\
2 & Does any nonce leak? & $s_1=c^{-1}(z-y)$\\
3 & Are rejection checks removed or changed? & ILWE regression on $cs$\\
4 & Is accepted $z$ non-uniform? & Fisher information bound\\
5 & What do rejected attempts reveal? & ILP on rejection bounds~\cite{rejected}\\
6 & Can a corrupt party bias its nonce share? & Checks always pass\\
7 & Are nonces tied to one group and session? & Nonce reuse\\
8 & Gaps in the attacker model? & Static-only proofs; $|S\setminus C|\ge2$\\
9 & How much do hints weaken the key? & Lattice-estimator bit loss\\
10 & Does code match the specification? & Float samplers, timing\\
\bottomrule
\end{tabularx}
\end{table}

\subsection{The attack methods in more detail}
\textbf{Least squares.} Each signature gives a pair $(c,b)$ with $b=-cs'+e$. Putting $N$ of these together gives a system
with more equations than unknowns. Its least-squares solution has an error per coefficient with variance about
$\sigma_e^2/(N\tau)$. For uniform noise on $[-\gamma_2,\gamma_2]$, $\sigma_e^2\approx\gamma_2^2/3$. Rounding works once
$\sigma_e/\sqrt{N\tau}<1/2$, which gives (\ref{eq:ilwe}).

\textbf{Bounded noise.} Under BCC the noise is known to lie in a fixed interval. So each sample gives two inequalities and
rules out every guess where $|b+cs'|$ is too large. The set of possible secrets shrinks faster than least squares would
suggest. We will build this as an ILP, following~\cite{rejected}, and also as a faster method that removes impossible
values step by step.

\textbf{Robust regression.} When only some samples carry information, Huber and Cauchy losses give less weight to the
useless ones. They work where least squares fails~\cite{concealed}.

\textbf{Fisher information.} If accepted responses follow a distribution that is shifted by $cs$, the Fisher information
$I$ per signature sets a limit (the Cram\'er--Rao bound) on how well any attack can estimate $s$. Recovering the key then
needs about $4/(I\tau)$ signatures~\cite{talus}. We will compute $I$ numerically from each scheme's exact acceptance region.

\textbf{Hints.} Partial knowledge of the secret, such as Mithril's share subsets, is written as extra linear equations
and given to a lattice estimator to measure how many bits of security are lost.

\textbf{How a simulator works (TALUS example).} For each signature the simulator: (1) draws the total nonce $y$ from the
scheme's distribution and keeps only nonces that meet BCC; (2) computes $w=Ay$ and $w_1=\HB(w)$; (3) makes $c$ from a
random message and $w_1$; (4) computes $z=y+cs_1$ and applies the first check; (5) outputs the public record $(c,z,h)$, plus
anything else the scheme broadcasts. It also stores the true answer, which the attack never sees. A second version also
outputs the internal shares of the corrupt parties. For schemes with aborts, step (4) outputs a record marked as rejected,
containing exactly the fields the specification says are revealed, and nothing more.

\textbf{Implementation.} Being correct and repeatable matters more than being fast. Arithmetic uses the NTT with
$q=8380417$ and a 512th root of unity, as in FIPS~204. The samplers (ExpandA, ExpandS, ExpandMask, SampleInBall) use
SHAKE128 and SHAKE256. We write prototypes in Python with vectorised maths, and move only the loops over $10^8$ to $10^9$
signatures to compiled code. Every experiment uses a fixed random seed and is logged, so any reported recovery can be
repeated.

\subsection{Method}
Step 1 repeats Niot's ILWE attack~\cite{niot} on TALUS-style data at reduced size and checks (\ref{eq:ilwe}). Step 2 (P1)
replaces least squares with bounded-noise and ILP methods and plots the success rate against the number of signatures.
Steps 3 and 4 (P2, P3) run the same methods on simulators of Mithril, Quorus and Trilithium. A negative result is stated
as ``no link to the key found within a given number of samples''. Step 5 (P4) is direct algebra on the published
protocol. Once full packages appear (from March 2027), we rerun the harness on the authors' own code.

Every experiment reports the number of signatures used, how much of the key was recovered, the running time and the
parameter set. We check our work in stages. The building blocks must pass the FIPS~204 test vectors. Each simulator must
be broken in a deliberately leaky version (for example one that publishes $Ay$), to show that the attacks actually work.
Negative results are reported together with the sample budget and the attacks that were tried.

\textbf{How we measure results.} We report (i) the smallest number of samples at which the whole secret is recovered at
least 90\% of the time over repeated runs; (ii) the fraction of coefficients recovered as the number of samples grows;
(iii) for statistical tests, the test statistic or $p$-value against the assumption that the data does not depend on the
key; and (iv) CPU time and memory. Every recovery is checked against the public key, so anyone can confirm a claimed break
by recomputing $t$. This prevents us from reporting a weak correlation as a break.

\textbf{What could make our results wrong.} Inside the project: bugs in the building blocks (handled by test vectors), and
attacks tuned too closely to our made-up keys (handled by testing many random keys). Outside the project: our simulators
might not match the real protocols, and parameters may change between the preview writeups and the final packages. In how
we measure: a scheme might leak in a way that one particular attack misses. So we run several attacks on every kind of
transcript, and we phrase negative results only as limits on what we tried.

\begin{table}[t]
\caption{Sample estimate (\ref{eq:ilwe}) per parameter set}\label{tab:ilwe}
\centering\footnotesize
\begin{tabular}{@{}lrrrr@{}}
\toprule
Set & $\gamma_2$ & $\tau$ & $4\gamma_2^2/(3\tau)$ & $\log_2$\\
\midrule
ML-DSA-44 & 95{,}232 & 39 & $3.1\times10^8$ & 28.2\\
ML-DSA-65 & 261{,}888 & 49 & $1.9\times10^9$ & 30.8\\
ML-DSA-87 & 261{,}888 & 60 & $1.5\times10^9$ & 30.5\\
\bottomrule
\end{tabular}
\end{table}

\subsection{First result}\label{sec:prelim}
Table~\ref{tab:ilwe} evaluates (\ref{eq:ilwe}). These $2^{28}$ to $2^{31}$ signatures are $2^{12}$ to $2^{15}$ times the TALUS
cap. A first version of the harness, with 18 unit tests, simulates $b=-cs+e$ with $e=\LB(w)$ for ML-DSA-44 and recovers $s$
by least squares. The measured error matches $\gamma_2/\sqrt{3\tau N}$ to within 3\% for $N$ from 5{,}000 to 160{,}000. If we
extend this trend, rounding all 1{,}024 coefficients of $s_2$ correctly needs about $3.6\times10^9$ signatures (for a 50\%
success rate). That is about 12 times the value from (\ref{eq:ilwe}). Our model leaves out $t_0$, bounded noise and lattice
reduction, so the formula should be read only as a rough order of magnitude.

\subsection{Phase 2 plan}
Table~\ref{tab:plan} lists the stages. Each one has a clear pass or fail test. If the building blocks or the leaky-version
tests fail, we stop and fix them first, because nothing built on top could be trusted.

\begin{table}[t]
\caption{Phase 2 plan}\label{tab:plan}
\centering\footnotesize
\begin{tabularx}{\columnwidth}{@{}lXX@{}}
\toprule
Stage & Work item & Output\\
\midrule
1 & Building blocks, transcript format; repeat ILWE check & Harness v0.1 passing test vectors\\
2 & TALUS cap: bounded-noise and ILP attacks (P1) & Success vs.\ sample curves\\
3 & Mithril accepted $z$, hint accounting (P2) & Fisher bound, bit-loss table\\
4 & Quorus and Trilithium rejected attempts (P3) & Statistical test results\\
5 & Kao P2 and DKG claims (P4) & Confirmed or refuted note\\
6 & Rerun on the authors' code & Updated results\\
\bottomrule
\end{tabularx}
\end{table}

\textbf{Is this doable?} Each problem comes down to a limited computation. P1 is linear algebra on simulated data. P2 is a
numerical integral plus one estimator call. P3 is a two-sample statistical test. P4 is algebra on a small example. None of
them needs private code or solving hard lattice problems at full size, because the attacks use leaks in the protocol
rather than breaking MLWE itself.

\textbf{Milestones.} M1: the building blocks match the FIPS~204 test vectors for all three versions. M2: a deliberately
leaky scheme (one that publishes $As_1$) is broken from one signature by Gaussian elimination. M3: the TALUS simulator
without the second check is broken by least squares at the number of signatures predicted by $\gamma_2/\sqrt{3\tau N}$,
and the factor-12 gap from Section~\ref{sec:prelim} is explained. M4: the bounded-noise attack beats least squares, or we
explain why it does not. M5 and M6: the Mithril, Quorus and Trilithium simulators follow their specifications. M7: for
each scheme we report either an attack or a bound.

\textbf{Possible outcomes and risks.} A positive result is an attack or a measured margin. A negative result is a statement
like ``no leak found with $q$ samples'', which is still useful evidence. The risks are that experiments with billions of
samples need compiled code, that our simulators may not match the real protocols if we have misread papers we have not
yet read in full, and that Mithril, whose team includes the author of~\cite{niot}, is probably already protected against
these kinds of attacks.

\textbf{What the project will deliver.} First, a reusable harness with tested ML-DSA building blocks and a simple
simulator interface, so a new threshold design can be checked in hours. Second, measured margins for the four proposals,
in signatures and in bits, so they can be compared on the same scale. Third, a confirmed or refuted statement about
Kao's profile P2, which matters because TALUS is built on it. Several of the failures we found came from copying ideas
from discrete-log schemes, so a tool that checks for these lattice-specific mistakes is useful beyond these four schemes.

\textbf{Limitations of this survey.} We have not yet read the full Mithril, Quorus and Trilithium papers from start to
finish. So the weak points listed for them are hypotheses that we still need to check against the original texts.
Several sources are preview writeups that may change before the deadline, and two of the cited ePrints have been
withdrawn. The TALUS numbers in~\cite{niot} differ slightly between the abstract and the table; we use the table. Phase 2
therefore starts by reading these papers in full.

\textbf{Ethics.} We only study public specifications that NIST has invited people to examine. If we find a problem, we
will tell the scheme's authors before making it public. We will clearly separate confirmed breaks, measured margins and
negative results, and say which results come from our simulators rather than the authors' code.

% =====================================================================
\section{Conclusion}
Threshold ML-DSA is moving from research towards standardisation. The published break of early TALUS versions shows how
easily splitting the key can leak it, either through commitments or through a missing rejection check. Our survey of
twenty-four sources finds three schemes with no public attacks, mainly because nobody has tested them, and one scheme
whose security now depends on a signature cap with an uncertain margin. It also finds that the attack tools (ILWE, ILP,
robust regression) are now good enough to test all four. We stated four open problems and proposed a leakage-testing
harness, and a first experiment confirmed how the least-squares error behaves. In Phase 2 we will complete the harness,
repeat the TALUS attack, and run the attacks on the other three schemes.

% =====================================================================
\begin{thebibliography}{24}\footnotesize
\bibitem{fips204} NIST, ``FIPS 204: Module-lattice-based digital signature standard,'' Aug. 2024.
\bibitem{dilithium} L. Ducas, E. Kiltz, T. Lepoint, V. Lyubashevsky, P. Schwabe, G. Seiler, and D. Stehl\'e, ``CRYSTALS-Dilithium: A lattice-based digital signature scheme,'' \emph{IACR TCHES}, vol. 2018, no. 1, pp. 238--268, 2018.
\bibitem{ir8214c} NIST, ``Multi-party threshold schemes: First call for submissions,'' NIST IR 8214C, 2026.
\bibitem{mithril} S. Celi, R. del Pino, T. Espitau, G. Niot, and T. Prest, ``Efficient threshold ML-DSA,'' in \emph{Proc. USENIX Security}, 2026 (IACR ePrint 2026/013).
\bibitem{reloaded} G. Borin, S. Celi, R. del Pino, T. Espitau, G. Niot, and T. Prest, ``Threshold signatures reloaded: ML-DSA and enhanced Raccoon with identifiable aborts,'' IACR ePrint 2025/1166, 2025 (withdrawn).
\bibitem{mithrilposter} S. Celi \emph{et al.}, ``Poster: Efficient threshold ML-DSA up to 6 parties,'' in \emph{Proc. ACM CCS}, 2025, doi:10.1145/3719027.3760739.
\bibitem{mithrilpw} S. Celi, G. Delerue, R. del Pino, T. Espitau, G. Niot, and T. Prest, ``Mithril: Efficient threshold ML-DSA from secret sharing with short shares,'' NIST MPTC preview writeup v1.0, Jan. 2026.
\bibitem{quorus} A. Bienstock, L. de Castro, D. Escudero, A. Polychroniadou, and A. Takahashi, ``Quorus: Efficient, scalable threshold ML-DSA signatures from MPC,'' in \emph{Proc. USENIX Security}, 2026 (IACR ePrint 2025/1163).
\bibitem{trilithium} A. Dufka, S. Kravt\v{s}enko, P. Laud, and N. Snetkov, ``Trilithium: Efficient and universally composable distributed ML-DSA signing,'' IACR ePrint 2025/675, 2025.
\bibitem{talus} L. Kao and R. Chang, ``TALUS: FIPS-204-exact threshold ML-DSA via boundary clearance,'' arXiv:2603.22109, 2026.
\bibitem{taluspw} L. Kao and R. Chang, ``TALUS (preview writeup v0.22),'' NIST MPTC First Call, Aug. 2026.
\bibitem{kao} L. Kao, ``FIPS 204-compatible threshold ML-DSA via Shamir nonce DKG,'' arXiv:2601.20917, 2026.
\bibitem{niot} G. Niot, ``Key-recovery attacks on TALUS: A cryptanalytic note,'' IACR ePrint 2026/1386, Jul. 2026.
\bibitem{raccoon} R. del Pino, S. Katsumata, M. Maller, F. Mouhartem, T. Prest, and M.-J. Saarinen, ``Threshold Raccoon: Practical threshold signatures from standard lattice assumptions,'' in \emph{Proc. EUROCRYPT}, LNCS 14652, 2024.
\bibitem{dpia} R. del Pino, T. Espitau, G. Niot, and T. Prest, ``Simple and efficient lattice threshold signatures with identifiable aborts,'' IACR ePrint 2025/871, 2025 (withdrawn).
\bibitem{gks} K. D. G\"ur, J. Katz, and T. Silde, ``Two-round threshold lattice-based signatures from threshold homomorphic encryption,'' in \emph{Proc. PQCrypto}, LNCS 14772, 2024.
\bibitem{ilwe} J. Bootle, C. Delaplace, T. Espitau, P.-A. Fouque, and M. Tibouchi, ``LWE without modular reduction and improved side-channel attacks against BLISS,'' in \emph{Proc. ASIACRYPT}, LNCS 11272, pp. 494--524, 2018.
\bibitem{rejected} Y. Zhou, W. Wang, Y. Sun, and Y. Yu, ``Rejected signatures' challenges pose new challenges: Key recovery of CRYSTALS-Dilithium via side-channel attacks,'' \emph{IACR TCHES}, 2025.
\bibitem{concealed} S. Damm \emph{et al.}, ``Solving concealed ILWE and its application for breaking masked Dilithium,'' in \emph{Proc. ASIACRYPT}, 2025.
\bibitem{yates} K. Yates, A. Pierrottet, A. Al Mamun, R. Cartor, M. Chowdhury, and S. Gao, ``Security analysis of integer learning with errors with rejection sampling,'' arXiv:2512.08172, 2025.
\bibitem{shamir} A. Shamir, ``How to share a secret,'' \emph{Commun. ACM}, vol. 22, no. 11, pp. 612--613, 1979.
\bibitem{feldman} P. Feldman, ``A practical scheme for non-interactive verifiable secret sharing,'' in \emph{Proc. IEEE FOCS}, 1987, pp. 427--438.
\bibitem{frost} C. Komlo and I. Goldberg, ``FROST: Flexible round-optimized Schnorr threshold signatures,'' in \emph{Proc. SAC 2020}, LNCS 12804, 2021, pp. 34--65.
\bibitem{finally} R. del Pino and G. Niot, ``Finally! A compact lattice-based threshold signature,'' in \emph{Proc. PKC}, LNCS 15676, 2025, pp. 169--199.
\end{thebibliography}

\end{document}
